\subsection{Worked solutions on the ten soils (third scenario)}
The same questions are repeated for different working conditions, to show how the answer moves with them.

\exbox{Ten-soil example 13c: Coarse loam (S01)}{Coarse loam ($K_s=13.0$ mm\,h$^{-1}$): a 5.6-mm nozzle at 400 kPa works on a $12\times18$ m spacing. Find the discharge, the application rate and whether runoff will occur.}\sol \givenline{$K_s=13.0$ mm\,h$^{-1}$}\ans $h=350/9.81=40.8$ m; \[q=0.95\pi(40.8/2000)^2\sqrt{2\times9.81\times5.6}\times6\times10^4=39.7\ \text{L\,min$^{-1}$}\] \ans \[r=q\times60/(S_lS_m)=39.7\times60/216=11.0\ \text{mm\,h$^{-1}$}\] \ans The rate is below $K_s$, so no runoff. \finans

\exbox{Ten-soil example 13c: Coarse sand (S02)}{Coarse sand ($K_s=120.0$ mm\,h$^{-1}$): a 5.6-mm nozzle at 400 kPa works on a $12\times18$ m spacing. Find the discharge, the application rate and whether runoff will occur.}\sol \givenline{$K_s=120.0$ mm\,h$^{-1}$}\ans $h=350/9.81=40.8$ m; \[q=0.95\pi(40.8/2000)^2\sqrt{2\times9.81\times5.6}\times6\times10^4=39.7\ \text{L\,min$^{-1}$}\] \ans \[r=q\times60/(S_lS_m)=39.7\times60/216=11.0\ \text{mm\,h$^{-1}$}\] \ans The rate is below $K_s$, so no runoff. \finans

\exbox{Ten-soil example 13c: Sand (S03)}{Sand ($K_s=210.0$ mm\,h$^{-1}$): a 5.6-mm nozzle at 400 kPa works on a $12\times18$ m spacing. Find the discharge, the application rate and whether runoff will occur.}\sol \givenline{$K_s=210.0$ mm\,h$^{-1}$}\ans $h=350/9.81=40.8$ m; \[q=0.95\pi(40.8/2000)^2\sqrt{2\times9.81\times5.6}\times6\times10^4=39.7\ \text{L\,min$^{-1}$}\] \ans \[r=q\times60/(S_lS_m)=39.7\times60/216=11.0\ \text{mm\,h$^{-1}$}\] \ans The rate is below $K_s$, so no runoff. \finans

\exbox{Ten-soil example 13c: Loamy sand (S04)}{Loamy sand ($K_s=61.0$ mm\,h$^{-1}$): a 5.6-mm nozzle at 400 kPa works on a $12\times18$ m spacing. Find the discharge, the application rate and whether runoff will occur.}\sol \givenline{$K_s=61.0$ mm\,h$^{-1}$}\ans $h=350/9.81=40.8$ m; \[q=0.95\pi(40.8/2000)^2\sqrt{2\times9.81\times5.6}\times6\times10^4=39.7\ \text{L\,min$^{-1}$}\] \ans \[r=q\times60/(S_lS_m)=39.7\times60/216=11.0\ \text{mm\,h$^{-1}$}\] \ans The rate is below $K_s$, so no runoff. \finans

\exbox{Ten-soil example 13c: Sandy loam (S05)}{Sandy loam ($K_s=26.0$ mm\,h$^{-1}$): a 5.6-mm nozzle at 400 kPa works on a $12\times18$ m spacing. Find the discharge, the application rate and whether runoff will occur.}\sol \givenline{$K_s=26.0$ mm\,h$^{-1}$}\ans $h=350/9.81=40.8$ m; \[q=0.95\pi(40.8/2000)^2\sqrt{2\times9.81\times5.6}\times6\times10^4=39.7\ \text{L\,min$^{-1}$}\] \ans \[r=q\times60/(S_lS_m)=39.7\times60/216=11.0\ \text{mm\,h$^{-1}$}\] \ans The rate is below $K_s$, so no runoff. \finans

\exbox{Ten-soil example 13c: Loam (S06)}{Loam ($K_s=13.0$ mm\,h$^{-1}$): a 5.6-mm nozzle at 400 kPa works on a $12\times18$ m spacing. Find the discharge, the application rate and whether runoff will occur.}\sol \givenline{$K_s=13.0$ mm\,h$^{-1}$}\ans $h=350/9.81=40.8$ m; \[q=0.95\pi(40.8/2000)^2\sqrt{2\times9.81\times5.6}\times6\times10^4=39.7\ \text{L\,min$^{-1}$}\] \ans \[r=q\times60/(S_lS_m)=39.7\times60/216=11.0\ \text{mm\,h$^{-1}$}\] \ans The rate is below $K_s$, so no runoff. \finans

\exbox{Ten-soil example 13c: Silt loam (S07)}{Silt loam ($K_s=6.8$ mm\,h$^{-1}$): a 5.6-mm nozzle at 400 kPa works on a $12\times18$ m spacing. Find the discharge, the application rate and whether runoff will occur.}\sol \givenline{$K_s=6.8$ mm\,h$^{-1}$}\ans $h=350/9.81=40.8$ m; \[q=0.95\pi(40.8/2000)^2\sqrt{2\times9.81\times5.6}\times6\times10^4=39.7\ \text{L\,min$^{-1}$}\] \ans \[r=q\times60/(S_lS_m)=39.7\times60/216=11.0\ \text{mm\,h$^{-1}$}\] \ans The rate exceeds $K_s$: ponding after $8.7$ h; sets longer than that run off, so use a smaller nozzle or a wider spacing. \finans

\exbox{Ten-soil example 13c: Sandy clay loam (S08)}{Sandy clay loam ($K_s=4.3$ mm\,h$^{-1}$): a 5.6-mm nozzle at 400 kPa works on a $12\times18$ m spacing. Find the discharge, the application rate and whether runoff will occur.}\sol \givenline{$K_s=4.3$ mm\,h$^{-1}$}\ans $h=350/9.81=40.8$ m; \[q=0.95\pi(40.8/2000)^2\sqrt{2\times9.81\times5.6}\times6\times10^4=39.7\ \text{L\,min$^{-1}$}\] \ans \[r=q\times60/(S_lS_m)=39.7\times60/216=11.0\ \text{mm\,h$^{-1}$}\] \ans The rate exceeds $K_s$: ponding after $3.8$ h; sets longer than that run off, so use a smaller nozzle or a wider spacing. \finans

\exbox{Ten-soil example 13c: Clay loam (S09)}{Clay loam ($K_s=2.3$ mm\,h$^{-1}$): a 5.6-mm nozzle at 400 kPa works on a $12\times18$ m spacing. Find the discharge, the application rate and whether runoff will occur.}\sol \givenline{$K_s=2.3$ mm\,h$^{-1}$}\ans $h=350/9.81=40.8$ m; \[q=0.95\pi(40.8/2000)^2\sqrt{2\times9.81\times5.6}\times6\times10^4=39.7\ \text{L\,min$^{-1}$}\] \ans \[r=q\times60/(S_lS_m)=39.7\times60/216=11.0\ \text{mm\,h$^{-1}$}\] \ans The rate exceeds $K_s$: ponding after $1.5$ h; sets longer than that run off, so use a smaller nozzle or a wider spacing. \finans

\exbox{Ten-soil example 13c: Silty clay loam (S10)}{Silty clay loam ($K_s=1.5$ mm\,h$^{-1}$): a 5.6-mm nozzle at 400 kPa works on a $12\times18$ m spacing. Find the discharge, the application rate and whether runoff will occur.}\sol \givenline{$K_s=1.5$ mm\,h$^{-1}$}\ans $h=350/9.81=40.8$ m; \[q=0.95\pi(40.8/2000)^2\sqrt{2\times9.81\times5.6}\times6\times10^4=39.7\ \text{L\,min$^{-1}$}\] \ans \[r=q\times60/(S_lS_m)=39.7\times60/216=11.0\ \text{mm\,h$^{-1}$}\] \ans The rate exceeds $K_s$: ponding after $1.2$ h; sets longer than that run off, so use a smaller nozzle or a wider spacing. \finans

